\section{Równania sterowania}
\label{sec:rownania}

\subsection{Zasada budowy}
Każda ze zmiennych stanu $M_1$, $M_2$, $M_3$ jest przechowywana
w przerzutniku RS. Dla przerzutnika o wejściach $S_x$ (Set) i $R_x$
(Reset) obowiązuje
\begin{equation}
    M_x^{\text{next}} = \bigl(S_x + M_x\bigr)\,\overline{R_x}.
    \label{eq:rs}
\end{equation}
Warunki $S_x$ i $R_x$ wyprowadzono z tabeli przejść
(tab.~\ref{tab:przejscia}). Dobrano je tak, aby były rozłączne, tzn.
$S_x R_x \equiv 0$.

\subsection{Telefon ($M_1$)}
Telefon zmienia stan wyłącznie po kliknięciu D. Jeżeli był wyłączony,
zostaje włączony, w przeciwnym razie wyłączony:
\begin{align}
    S_1 &= M_6\,\overline{M_1}, \label{eq:s1}\\
    R_1 &= M_6\,M_1. \label{eq:r1}
\end{align}
Ponieważ $M_6$ trwa jeden cykl, a warunki $S_1$ i $R_1$ zależą od
$M_1$ w przeciwny sposób, zmiana stanu następuje dokładnie raz na
kliknięcie. Przy impulsie dłuższym niż jeden cykl układ w sposób
ciągły przełączałby $M_1$.

\subsection{Ekran ($M_2$)}
\begin{align}
    S_2 &= M_4\,M_1\,\overline{M_3}\,\overline{M_2} \;+\; M_6\,\overline{M_1}, \label{eq:s2}\\
    R_2 &= M_4\,M_1\,\overline{M_3}\,M_2 \;+\; M_6\,M_1. \label{eq:r2}
\end{align}
Pierwsze składniki realizują przełączanie ekranu kliknięciem K
(wymaganie R2): gdy $M_2=0$ aktywne jest Set, gdy $M_2=1$ -- Reset.
Czynnik $\overline{M_3}$ powoduje, że przy otwartym menu kliknięcie K
nie zmienia $M_2$, a więc po anulowaniu menu ekran wraca do
poprzedniego stanu (R4). Składnik $M_6\,\overline{M_1}$ włącza ekran
przy uruchamianiu telefonu ($S_0 \to S_2$), a $M_6\,M_1$ zeruje go
przy wyłączaniu.

\subsection{Menu ($M_3$)}
\begin{align}
    S_3 &= M_5\,M_1, \label{eq:s3}\\
    R_3 &= M_4\,M_3 \;+\; M_6. \label{eq:r3}
\end{align}
Kliknięcie Ś otwiera menu tylko przy włączonym telefonie
(R1, R3). Menu zamykane jest przez K (anulowanie, R4) lub przez D
(wyłączenie telefonu, R6). Brak składnika $M_6$ w $R_3$ skutkowałby
pozostaniem menu otwartego po wyłączeniu telefonu.

\subsection{Uproszczenie i wybór postaci równań}
W równaniu~\eqref{eq:s2} czynnik $\overline{M_2}$ w pierwszym
składniku można by pominąć, jeśli przyjąć, że w bloku RS wejście
Reset ma priorytet nad Set. Wtedy przy $M_2=1$ aktywne byłyby oba
wejścia, a o stanie decydowałaby zasada ,,Reset wygrywa''. Nie
zastosowano tego uproszczenia,
ponieważ postać z jawnymi czynnikami zapewnia rozłączność $S_2$
i $R_2$, a poprawność układu nie zależy wówczas od priorytetu wejść
bloku RS. Z tego samego powodu pozostałe pary $S_x$, $R_x$ są
rozłączne: $S_1 R_1 = 0$ (czynniki $\overline{M_1}$ i $M_1$),
a zdarzenia $M_4$, $M_5$, $M_6$ nigdy nie występują jednocześnie.

Równanie~\eqref{eq:q1} można by uprościć do $Q_1 = M_2 + M_3$, gdyż
w punkcie~\ref{sec:niedozwolone} pokazano, że $M_1 = 0$ implikuje
$M_2 = M_3 = 0$. Czynnik $M_1$ został jednak zachowany jako
zabezpieczenie (por. R7).

\subsection{Weryfikacja równań względem tabeli przejść}
Tabela~\ref{tab:weryfikacja} zawiera wynik podstawienia kodów
stanów i zdarzeń do równań~\eqref{eq:s1}--\eqref{eq:r3}. W każdym
wierszu stan następny wynikający z równań jest zgodny z tabelą
przejść~\ref{tab:przejscia}.

\begin{table}[H]
\centering
\caption{Weryfikacja równań sterowania dla wszystkich stanów i zdarzeń.}
\label{tab:weryfikacja}
\begin{tabular}{lll|c|l}
\toprule
Stan & Zdarzenie & Aktywne wejścia RS & $M_1M_2M_3$ & Stan następny \\
\midrule
$S_0$    & K, Ś & --- & 000 & $S_0$ \\
$S_0$    & D & $S_1$, $S_2$, $R_3$ & 110 & $S_2$ \\
$S_1$    & K & $S_2$ & 110 & $S_2$ \\
$S_1$    & Ś & $S_3$ & 101 & $S_{3a}$ \\
$S_1$    & D & $R_1$, $R_2$, $R_3$ & 000 & $S_0$ \\
$S_2$    & K & $R_2$ & 100 & $S_1$ \\
$S_2$    & Ś & $S_3$ & 111 & $S_{3b}$ \\
$S_2$    & D & $R_1$, $R_2$, $R_3$ & 000 & $S_0$ \\
$S_{3a}$ & K & $R_3$ & 100 & $S_1$ \\
$S_{3a}$ & Ś & $S_3$ & 101 & $S_{3a}$ \\
$S_{3a}$ & D & $R_1$, $R_2$, $R_3$ & 000 & $S_0$ \\
$S_{3b}$ & K & $R_3$ & 110 & $S_2$ \\
$S_{3b}$ & Ś & $S_3$ & 111 & $S_{3b}$ \\
$S_{3b}$ & D & $R_1$, $R_2$, $R_3$ & 000 & $S_0$ \\
\bottomrule
\end{tabular}
\end{table}

\subsection{Brak stanów niedozwolonych}
\label{sec:niedozwolone}
Wymaganie R7 wymaga wykluczenia stanów, w których $M_1=0$ przy
$M_2=1$ lub $M_3=1$. Wynika to z równań:
\begin{itemize}
    \item $M_1$ zmienia się wyłącznie w chwili $M_6=1$, a wtedy
          zachodzi: przy wyłączaniu ($M_6 M_1$) wejścia $R_2$ i $R_3$
          zerują $M_2$ i $M_3$, przy włączaniu ($M_6 \overline{M_1}$)
          $M_2$ zostaje ustawione, a $M_3$ zerowane przez $R_3$,
    \item $M_3$ może zostać ustawione tylko przy $M_1=1$ (równanie~\eqref{eq:s3}),
    \item $M_2$ może zostać ustawione tylko przy $M_1=1$ (przełączanie K)
          albo w chwili włączania telefonu.
\end{itemize}
Po wyłączeniu telefonu zawsze $M_2 = M_3 = 0$. Układ nie opuszcza więc
zbioru pięciu dozwolonych stanów z tabeli~\ref{tab:stany}.
